\documentclass[../Main.tex]{subfiles}

\begin{document}

\section{Reflexion} \label{sec:Reflexion}

\subsection{Reflexion des Arbeitsprozesses}

Dieses Projekt hat etwa ein Jahr unserer Zeit in Anspruch genommen und mehrere Phasen durchlaufen, die wir im Folgenden der Reihe nach aufgreifen.

Ganz am Anfang waren wir noch völlig unsicher und hatten ganz andere Vorstellungen davon, welches Produkt am Ende entstehen würde. Also stocherten wir in alle Bereiche hinein: Wir begannen mit Spielsystemen, von denen wir ohne Recherche annahmen, dass wir sie brauchen, überlegten, wie wir die Erkenntnisfrage beantworten wollen, und recherchierten, wie man recherchiert. Bis zur Zwischenpräsentation blieben uns noch vier Monate. Wir waren wenig zielorientiert und machten einfach, was uns gefiel oder das was wir noch nicht ausprobiert hatten. Das mag planlos wirken, rückblickend war es aber ein Auskundschaften dessen, was es gibt. Die Fallhöhe war zudem gering: Was auch immer wir machten, nichts musste verworfen werden, und die Arbeitsschritte hingen noch kaum voneinander ab. Dass das aufging, war zwar eher ein glücklicher Zufall als strukturiertes Denken unsererseits, doch so ist die Natur eines Projektanfangs. Unter Leuten, die Projekte machen, heisst es ja, sie hätten nie angefangen, hätten sie gewusst, wie viel Arbeit es wird. Könnten wir noch einmal beginnen, würden wir es genauso angehen, mit dem einzigen Unterschied, dass wir wüssten, dass wir am Anfang eine Narrenfrist haben.

Bis zur Zwischenpräsentation spielte die Erkenntnisfrage eine untergeordnete Rolle, denn die dafür relevanten Systeme beruhten auf einer technischen Basis, die wir noch nicht hatten: einem Spielercharakter mit Animationen und Steuerung, Gegnern und Code, der das Spiel überhaupt zum Starten bringt. Obwohl wir dabei die Recherche vernachlässigten, war das rückblickend die richtige Entscheidung. Zum einen machte uns dieser Teil am meisten Freude, zum anderen mussten wir uns zuerst eine Vorstellung von unserem eigenen Spiel bilden. Erst nachdem wir die ersten Systeme gebaut hatten, wussten wir, worauf wir in anderen Spielen achten mussten, und konnten erst recht eine gute Recherche machen. Bei einem nächsten Projekt würden wir die Reihenfolge davon abhängig machen, wie konkret unsere Vorstellungen sind und wie gross unser Kenntnisstand ist: Haben wir beides, recherchieren wir zuerst. Fehlt eines davon, sammeln wir zuerst eigene Erfahrungen -- also genau das, was wir diesmal eher unbeabsichtigt getan haben.

Gleichzeitig rückte die Zwischenpräsentation näher. Features, die wir zwar wollten, die für die Erkenntnisfrage aber wenig relevant waren, mussten wir schon jetzt streichen oder stark eingrenzen, und wir mussten priorisieren, was für ein Spiel essenziell ist. Dazu kam ein begründet geplanter Playtest, und die Checkliste aus dem Vertrag hatte noch viele offene Punkte. Wir haben alles erreicht, was wir vertraglich festgehalten hatten, auch wenn es sehr, sehr knapp wurde. Nicht das Fertigstellen, sondern die Qualitätssicherung hat uns einen Strich durch die Rechnung gemacht, weil wir möglichst wenige Kompromisse eingehen wollten. Im Nachhinein wäre es interessant gewesen zu erfahren, wie 80 oder 85 statt 100 Prozent unseres Anspruchs bewertet worden wären. Da die Zwischenbewertung nicht in die Note der Maturitätsarbeit einfliesst, ist das eine verpasste Chance. Andererseits konnten wir Erfahrung mit Crunchen sammeln, die uns am Ende ermöglicht hat, mit genügend Zeit und kaum Crunchen 95 bis 100 Prozent herauszuholen.

Zur Zwischenpräsentation erhielten wir eine Rückmeldung: Vieles war gut, einiges weniger. Die Kritik betraf weniger den Prototyp als alles rundherum, besonders die begleitenden Texte. Beim Umsetzen hatten wir zunächst Schwierigkeiten, denn hätten wir die bemängelten Punkte verinnerlicht gehabt, hätten wir sie von Anfang an richtig gemacht. Deshalb stellten wir unserer Betreuungsperson so lange Fragen, bis wir sicher waren, sie verstanden zu haben.

Nach der Zwischenpräsentation hatten wir wieder genug Zeit, und die Abgabe war weit genug entfernt, dass wir unsere Ambitionen beim Projektumfang vergrösserten. Gleichzeitig nahm das anfängliche Tempo ab, mit dem wir Neues produzierten. Zum einen war das Projekt inzwischen so gross, dass kleinere Neuerungen wenig Unterschied machten. Zum anderen war unser anfänglicher Eifer zwar weit davon entfernt, erloschen zu sein, doch das Projekt zu pflegen und Aufgaben zu erledigen, die uns weniger Spass machten und sich mittlerweile angehäuft hatten, setzte uns langsam zu.

mitlerweile waren prozesse auch so abhängig voneinander, dass wir in unseerer partnerarbeit nicht einfach tun und lassen konnten, was wir owllten und wir mussten uns aufeiandner abstimmen. dabei arbeiteten wir an möglichst grossen blöcken die unabhängig voneinander waren und stimmten uns ab auf ein datum, wenn sich unsee arbeiten kreutzten. in dieser zeit wurde fast die gesammte restliche recherche gemacht. dies lief ser gut wenn beide rechtzeitig ferig wurdne, was nicht immer passiert ist. als verbesserungsvorschlag noch einen dritten görsseren strang zu haben, der offen ist, wenn eine person durch eine andere blockiert ist.

nach ende der recherche wussten wir dann auch genau was für systeme noch weiter implementiert werden mussten. mit der deadline immer näher kommend haben wir dann auch mittlerweile das beantworten der forschungsfrage im sinn. hier fängt dann unsere wirklich gut geplanten teile an. weil wir jetzt uns bewusst sind, dass nicht mehr jede arbeit nicht verschwendet ist sondern wir jetzt schon ein fundament haben, dass in eine richtugn geht, arbeiten wir zielgerichtet für die erkenntnisfrage. dieser teil ist uns besodners gelungen, da wir auf hinweis unserer begleitperson einen plan bis zum ende gemacht haben. und dass dieser plan gut geworden ist, hat damit zu tun, dass wir schon bei der zwischenabgabe erfahrungen gesammelt hatten, wie man auf eine abgabe hinarbeitet und plant.

gegen ende der arbeit haben wir noch unsere playtest machen müssen, die recht gut gelaufen sind, da wir schon selber playtest gemacht hatten für die rechrche. und in letzten schlusspahse machte uns eher formalitäten sorgen. ich glaube diese zeit ist auch ziemlich gut gelaufen, nicht weil wir in dieser phase etwas spezielles gemacht haben sondern weil wir in der phase davor gut geplant haben und uns genug zeit für diese phase gelassen haben. des weiteren waren wir uns bewusst, dass obwohl es noch nach genug zeit aussieht es kaapp wird wenn wir uns nicht dran halten.

\begin{itemize}[leftmargin=1.5cm]
    \item \textcolor{red}{\textbf{1. Ausgangslage und Plan:} Erste Commits im Juni 2025; der Vertrag sah vor: Zwischenprodukt 10.6., Zwischenpräsentation 17.6., Analyse der Referenzspiele Anfang Juli, Spielsysteme Mitte August, Doku und Fragebögen Ende September, Playtests Mitte Oktober, Verbesserungen Ende Oktober.}
    \item \textcolor{red}{\textbf{2. Verlauf in Phasen} (aus Git und Journal): Hauptentwicklung des Prototyps Februar bis Juni 2026; Juni: README und Architekturdiagramm fürs Zwischenprodukt; 16.--20.7.: Kapitel 1--4 und die \game{NieR}-Analyse als Entwurf; Juli bis September kaum Arbeit am Spiel (Sommerpause?); Oktober: Level überarbeitet, Playtests vorbereitet und durchgeführt; November: Begleitbericht fertiggestellt.}
    \item \textcolor{red}{\textbf{3. Scope-Entscheide:} 20.5.: Vertrag bereinigt, Projektumfang reduziert, Erkenntnisfrage festgelegt -- was wurde gestrichen? Später: Bonus-Inhalte weggelassen, technische Dokumentation ausgelagert, README gestrichen, Recherche in Kapitel~\ref{sec:Leveldesign} auf Rat der Betreuungsperson auf das für die Erkenntnisfrage Nötige gekürzt.}
    \item \textcolor{red}{\textbf{4. Plan und tatsächlicher Verlauf:} Playtests 26.10.--2.11.\ statt Mitte Oktober, Verbesserungen erst im Bericht statt Ende Oktober. Warum?}
    \item \textcolor{red}{\textbf{5. Wichtige Wendepunkte} (Journal 24./25.10., 13~h): Obwohl der Prototyp technisch fertig war, machten die Kämpfe bei eigenen Tests wenig Spass. Ihr habt die Encounters gekürzt, die Gegner aggressiver gemacht und Rechercheergebnisse nachgeholt, die noch nicht umgesetzt waren. Eure eigene Erkenntnis dort: \enquote{Die technische Fertigstellung allein garantiert noch kein gutes Spielerlebnis} -- und weil ihr das Level so oft gespielt habt, braucht es externe Playtests. Das ist ein starkes Reflexionsbeispiel.}
    \item \textcolor{red}{\textbf{6. Rückmeldungen der Zwischenbewertung und was ihr daraus gemacht habt:} Zeitaufwand konsequenter dokumentieren, Reflexionen vertiefen, Entscheidungen mit Alternativen begründen, KI vollständig dokumentieren, Planung mit Arbeitspaketen, Bericht und Code verknüpfen.}
    \item \textcolor{red}{\textbf{7. Was gut funktioniert hat:} z.\,B. iteratives Vorgehen, Git und GitHub, Playtest-Vorgehen -- jeweils mit Grund.}
    \item \textcolor{red}{\textbf{8. Methodik im Rückblick:} kurz, Einzelheiten in den Limitationen (Abschnitt~\ref{sec:Limitationen}).}
    \item \textcolor{red}{\textbf{9. Was ihr anders machen würdet}, z.\,B. früher extern testen, Rechercheergebnisse von Anfang an mit dem Level abgleichen, Bericht laufend statt am Schluss schreiben.}
    \item \textcolor{red}{\textbf{10. Verweis aufs Arbeitsjournal}, wo der Verlauf im Detail dokumentiert ist.}
\end{itemize}

\subsection{Gruppenarbeit und Aufgabenteilung}

\begin{itemize}[leftmargin=1.5cm]
    \item \textcolor{red}{\textbf{1. Wie ihr die Arbeit aufgeteilt habt.} Laut Zwischenbewertung: Rodin stärker bei Animation, UI, Bericht und README; Zhile stärker bei Architektur, Bewegung, State Machines und Kampf-Grundlagen. Das deckt sich mit Git: Zhile vor allem \code{Entities/Player}, \code{Common} (Utility, Tools), \code{Core}, Kamera; Rodin vor allem \code{Interfaces} und Menüs sowie die TrenchBroom-Karten des Levels.}
    \item \textcolor{red}{\textbf{2. Tabelle Aufgabenteilung und Eigenleistung} für Prototyp und Bericht (Rodin / Zhile / gemeinsam). Zum Bericht: Kapitel 1--8 und die \game{NieR}-Analyse Rodin, die \game{Ninja-Gaiden}-Analyse Zhile, README und Architekturdiagramm (Juni) vor allem Zhile. Bitte prüfen und ergänzen (Playtests, Leveldesign, Audio, \ldots).}
    \item \textcolor{red}{\textbf{3. Zusammenarbeit und Kommunikation:} Wie habt ihr euch abgesprochen (Treffen, Discord, GitHub)? Wie habt ihr Entscheidungen getroffen?}
    \item \textcolor{red}{\textbf{4. Veränderung im Verlauf:} z.\,B. ab Oktober Rodin vor allem Level, Playtests und Bericht, Zhile weiter am Spiel (laut Git im Oktober sehr aktiv). Wie kam es dazu?}
    \item \textcolor{red}{\textbf{5. Fazit:} Was lief gut, was war schwierig, was würdet ihr in einer Gruppenarbeit anders machen? Hinweis: Commit-Zahlen sagen wenig über den Aufwand aus, deshalb nicht als Massstab verwenden.}
\end{itemize}

\subsection{Verwendung von KI} \label{sec:ReflexionKI}

\begin{itemize}[leftmargin=1.5cm]
    \item \textcolor{red}{\textbf{1. Welche Tools wofür:} Code (Treppen-Werkzeug und Animations-Hilfsskript im Frühling, Audiosystem im Oktober), Begleitbericht (Korrektur, Struktur, Rückmeldungen). Details in der KI-Deklaration (Abschnitt~\ref{sec:KIDeklaration}).}
    \item \textcolor{red}{\textbf{2. Was geholfen hat:} z.\,B. Widersprüche zwischen Kapiteln finden, unbelegte Aussagen markieren, \LaTeX-Formatierung, Fehlersuche im Code.}
    \item \textcolor{red}{\textbf{3. Grenzen und Probleme:} Die KI ging mehrmals von falschen Annahmen aus, die ihr korrigieren musstet, z.\,B. dass Themenfindung mit KI gemacht wurde, dass ein Tutorial geplant war, welche Räume Brückenräume sind, oder dass in den Protokollen ein anderer Archer gemeint war.}
    \item \textcolor{red}{\textbf{4. Wie ihr die Ergebnisse überprüft habt:} Abgleich mit Quellen, Playtest-Protokollen und dem Spiel selbst; Vorschläge nicht ungeprüft übernommen, einige bewusst verworfen.}
    \item \textcolor{red}{\textbf{5. Eigenanteil sichern:} Interpretation und Urteile selbst festgelegt; von der KI formulierte Stellen markiert und in eigenen Worten neu geschrieben; Regel \enquote{nicht wortwörtlich übernehmen} aus dem Merkblatt.}
    \item \textcolor{red}{\textbf{6. Fazit:} Würdet ihr KI wieder so einsetzen? Was würdet ihr anders machen, z.\,B. KI-Nutzung von Anfang an laufend im Journal dokumentieren?}
\end{itemize}

\subsection{Persönliche Erkenntnisse}

\subsubsection{Rodin Arslan}

\begin{itemize}[leftmargin=1.5cm]
    \item \textcolor{red}{\textbf{1. Fachlich:} z.\,B. Animation (u.\,a. Hängeanimation), UI und Menüs, Leveldesign mit TrenchBroom, Gegner- und Encounter-Design, Playtesting, \LaTeX\ und BibTeX.}
    \item \textcolor{red}{\textbf{2. Methodisch:} Spielanalyse mit dem MDA-Framework, Recherche und Zitieren, Playtests planen und auswerten, Projektplanung.}
    \item \textcolor{red}{\textbf{3. Persönlich:} z.\,B. dass ein technisch fertiges Spiel noch nicht gut spielbar ist (Journal 24./25.10.), Zeitmanagement in der Schlussphase, Umgang mit Rückmeldungen.}
    \item \textcolor{red}{\textbf{4. Was du mitnimmst} für die Zukunft.}
\end{itemize}

\subsubsection{Zhile Zhu}

\textcolor{red}{ZHILE DAS IST FÜR DICH}

\end{document}
